% --- Archon physics formalization source begin ---
% archon:physics
% archon:covers PhyXMiniProblems/problem_phyx_mini_0298.lean
% archon:source-report reports/phyx_mini/problem_phyx_mini_0298.source.json

\chapter{Physics problem phyx\_mini\_0298}
\label{ch:PhyXMiniProblems_problem_phyx_mini_0298}

\paragraph{Problem source.}
\#\# Physical scenario

In figure, a string, tied to a sinusoidal oscillator at \$P\$ and running over a support at \$Q\$, is stretched by a block of mass \$m\$. Separation \$L = 1.20\$ m, linear density \$\textbackslash{}mu = 1.6\$ g/m, and the oscillator frequency \$f = 120\$ Hz. The amplitude of the motion at \$P\$ is small enough for that point to be considered a node. A node also exists at \$Q\$.

\#\# Question

What mass \$m\$ allows the oscillator to set up the fourth harmonic on the string?

\#\# Answer choices

- A: A: 0.764 kg

- B: B: 0.786 kg

- C: C: 0.815 kg

- D: D: 0.846 kg

\#\# Figure caption (auxiliary; use the image as primary evidence)

The image depicts a physics diagram illustrating wave motion along a string. 

- **Objects and Scenes**: 
  - There is an oscillator on the left side, represented by a blue rectangular object with the label "Oscillator".
  - A green wavy line signifies the wave traveling along the string from point P to point Q.
  - Points P and Q are marked with black dots at each end of the string.

- **Relationships**:
  - The oscillator initiates the wave at point P.
  - The string between P and Q is of length L.
  - Point Q is connected to a blue square block labeled as "m", indicating a mass at the end of the string.

- **Text**:
  - The labels "Oscillator", "P", "Q", "L", and "m" are present, with arrows indicating the physical parameters involved. 

The diagram portrays the dynamics of wave propagation and tension in a string attached to an oscillator.

\#\# Dataset metadata

Resonance and Harmonics; Physical Model Grounding Reasoning, Multi-Formula Reasoning

\paragraph{Recorded answer/context.}
D: D: 0.846 kg

\paragraph{Figure/image path.}
/root/proposal\_for\_physic/science-mango/phyx\_mini\_run/phyx\_data/test\_image/298.png

\paragraph{Formalization target.}
create a compiling Lean file with sorry bodies at `PhyXMiniProblems/problem\_phyx\_mini\_0298.lean`. The Lean declarations must preserve the physical quantities, dimensions or dimensional roles, figure labels, governing-law hypotheses, and final relation expressed by this problem.
Use Mathlib/Physlib names found through LeanExplore where available. If a physics API is missing, introduce faithful local abstractions rather than scalar placeholder aliases.

\begin{theorem}[Physics formalization target]
\label{thm:physics:phyx_mini_0298:target}
\lean{PhyXMiniProblems.ProblemPhyXMini0298.problem_phyx_mini_0298}
\uses{def:physics:phyx-mini-0298:phyxminiproblems-problemphyxmini0298-massinkilograms, def:physics:phyx-mini-0298:phyxminiproblems-problemphyxmini0298-oscillatorstringsetup, def:physics:phyx-mini-0298:phyxminiproblems-problemphyxmini0298-matchesproblemandprimaryfigure, def:physics:phyx-mini-0298:phyxminiproblems-problemphyxmini0298-usesstandardearthgravity, def:physics:phyx-mini-0298:phyxminiproblems-problemphyxmini0298-haspositivephysicalparameters, def:physics:phyx-mini-0298:phyxminiproblems-problemphyxmini0298-satisfiesidealsupportstringandharmoniclaws, def:physics:phyx-mini-0298:phyxminiproblems-problemphyxmini0298-recordedanswerchoice, def:physics:phyx-mini-0298:phyxminiproblems-problemphyxmini0298-matchesanswerchoice}
The assigned autoformalize agent should translate this physics problem into Lean declarations in the covered file, with theorem and lemma proof bodies written as `by sorry`.
\end{theorem}
\begin{proof}
This is an autoformalization task, not a proof task. Produce statements that can later be proved by the physics prover without weakening the physical model.
\end{proof}

% --- Archon declaration topology begin ---
\section{Declaration topology}
\begin{definition}[Mass Quantity]
  \label{def:physics:phyx-mini-0298:phyxminiproblems-problemphyxmini0298-massquantity}
  \lean{PhyXMiniProblems.ProblemPhyXMini0298.MassQuantity}
  A nonnegative physical mass
\end{definition}
\begin{proof}
  This is the declared model definition of Mass Quantity.
\end{proof}
\begin{definition}[Length Quantity]
  \label{def:physics:phyx-mini-0298:phyxminiproblems-problemphyxmini0298-lengthquantity}
  \lean{PhyXMiniProblems.ProblemPhyXMini0298.LengthQuantity}
  A nonnegative physical length
\end{definition}
\begin{proof}
  This is the declared model definition of Length Quantity.
\end{proof}
\begin{definition}[Frequency Quantity]
  \label{def:physics:phyx-mini-0298:phyxminiproblems-problemphyxmini0298-frequencyquantity}
  \lean{PhyXMiniProblems.ProblemPhyXMini0298.FrequencyQuantity}
  A nonnegative ordinary frequency, carrying inverse-time dimension
\end{definition}
\begin{proof}
  This is the declared model definition of Frequency Quantity.
\end{proof}
\begin{definition}[Acceleration Quantity]
  \label{def:physics:phyx-mini-0298:phyxminiproblems-problemphyxmini0298-accelerationquantity}
  \lean{PhyXMiniProblems.ProblemPhyXMini0298.AccelerationQuantity}
  A nonnegative acceleration magnitude
\end{definition}
\begin{proof}
  This is the declared model definition of Acceleration Quantity.
\end{proof}
\begin{definition}[Tension Quantity]
  \label{def:physics:phyx-mini-0298:phyxminiproblems-problemphyxmini0298-tensionquantity}
  \lean{PhyXMiniProblems.ProblemPhyXMini0298.TensionQuantity}
  A nonnegative string-tension magnitude, carrying force dimension
\end{definition}
\begin{proof}
  This is the declared model definition of Tension Quantity.
\end{proof}
\begin{definition}[Linear Mass Density Quantity]
  \label{def:physics:phyx-mini-0298:phyxminiproblems-problemphyxmini0298-linearmassdensityquantity}
  \lean{PhyXMiniProblems.ProblemPhyXMini0298.LinearMassDensityQuantity}
  A nonnegative mass per unit length of string
\end{definition}
\begin{proof}
  This is the declared model definition of Linear Mass Density Quantity.
\end{proof}
\begin{definition}[Speed Quantity]
  \label{def:physics:phyx-mini-0298:phyxminiproblems-problemphyxmini0298-speedquantity}
  \lean{PhyXMiniProblems.ProblemPhyXMini0298.SpeedQuantity}
  A unit-independent, nonnegative propagation speed
\end{definition}
\begin{proof}
  This is the declared model definition of Speed Quantity.
\end{proof}
\begin{definition}[quantity Readout]
  \label{def:physics:phyx-mini-0298:phyxminiproblems-problemphyxmini0298-quantityreadout}
  \lean{PhyXMiniProblems.ProblemPhyXMini0298.quantityReadout}
  Read a nonnegative dimensionful quantity in a coherent unit system
\end{definition}
\begin{proof}
  This is the declared model definition of quantity Readout.
\end{proof}
\begin{definition}[mass In Kilograms]
  \label{def:physics:phyx-mini-0298:phyxminiproblems-problemphyxmini0298-massinkilograms}
  \lean{PhyXMiniProblems.ProblemPhyXMini0298.massInKilograms}
  \uses{def:physics:phyx-mini-0298:phyxminiproblems-problemphyxmini0298-massquantity, def:physics:phyx-mini-0298:phyxminiproblems-problemphyxmini0298-quantityreadout}
  Kilogram readout of a physical mass
\end{definition}
\begin{proof}
  This is the declared model definition of mass In Kilograms.
\end{proof}
\begin{definition}[length In Meters]
  \label{def:physics:phyx-mini-0298:phyxminiproblems-problemphyxmini0298-lengthinmeters}
  \lean{PhyXMiniProblems.ProblemPhyXMini0298.lengthInMeters}
  \uses{def:physics:phyx-mini-0298:phyxminiproblems-problemphyxmini0298-lengthquantity, def:physics:phyx-mini-0298:phyxminiproblems-problemphyxmini0298-quantityreadout}
  Metre readout of the separation between P and Q
\end{definition}
\begin{proof}
  This is the declared model definition of length In Meters.
\end{proof}
\begin{definition}[frequency In Hertz]
  \label{def:physics:phyx-mini-0298:phyxminiproblems-problemphyxmini0298-frequencyinhertz}
  \lean{PhyXMiniProblems.ProblemPhyXMini0298.frequencyInHertz}
  \uses{def:physics:phyx-mini-0298:phyxminiproblems-problemphyxmini0298-frequencyquantity, def:physics:phyx-mini-0298:phyxminiproblems-problemphyxmini0298-quantityreadout}
  Hertz readout of the oscillator frequency
\end{definition}
\begin{proof}
  This is the declared model definition of frequency In Hertz.
\end{proof}
\begin{definition}[acceleration In Meters Per Second Squared]
  \label{def:physics:phyx-mini-0298:phyxminiproblems-problemphyxmini0298-accelerationinmeterspersecondsquared}
  \lean{PhyXMiniProblems.ProblemPhyXMini0298.accelerationInMetersPerSecondSquared}
  \uses{def:physics:phyx-mini-0298:phyxminiproblems-problemphyxmini0298-accelerationquantity, def:physics:phyx-mini-0298:phyxminiproblems-problemphyxmini0298-quantityreadout}
  Metres-per-second-squared readout of an acceleration
\end{definition}
\begin{proof}
  This is the declared model definition of acceleration In Meters Per Second Squared.
\end{proof}
\begin{definition}[tension In Newtons]
  \label{def:physics:phyx-mini-0298:phyxminiproblems-problemphyxmini0298-tensioninnewtons}
  \lean{PhyXMiniProblems.ProblemPhyXMini0298.tensionInNewtons}
  \uses{def:physics:phyx-mini-0298:phyxminiproblems-problemphyxmini0298-tensionquantity, def:physics:phyx-mini-0298:phyxminiproblems-problemphyxmini0298-quantityreadout}
  Newton readout of a string tension
\end{definition}
\begin{proof}
  This is the declared model definition of tension In Newtons.
\end{proof}
\begin{definition}[linear Mass Density In Kilograms Per Meter]
  \label{def:physics:phyx-mini-0298:phyxminiproblems-problemphyxmini0298-linearmassdensityinkilogramspermeter}
  \lean{PhyXMiniProblems.ProblemPhyXMini0298.linearMassDensityInKilogramsPerMeter}
  \uses{def:physics:phyx-mini-0298:phyxminiproblems-problemphyxmini0298-linearmassdensityquantity, def:physics:phyx-mini-0298:phyxminiproblems-problemphyxmini0298-quantityreadout}
  Kilograms-per-metre readout of a linear mass density
\end{definition}
\begin{proof}
  This is the declared model definition of linear Mass Density In Kilograms Per Meter.
\end{proof}
\begin{definition}[speed In Meters Per Second]
  \label{def:physics:phyx-mini-0298:phyxminiproblems-problemphyxmini0298-speedinmeterspersecond}
  \lean{PhyXMiniProblems.ProblemPhyXMini0298.speedInMetersPerSecond}
  \uses{def:physics:phyx-mini-0298:phyxminiproblems-problemphyxmini0298-speedquantity, def:physics:phyx-mini-0298:phyxminiproblems-problemphyxmini0298-quantityreadout}
  Metres-per-second readout of a transverse wave speed
\end{definition}
\begin{proof}
  This is the declared model definition of speed In Meters Per Second.
\end{proof}
\begin{definition}[Figure Point]
  \label{def:physics:phyx-mini-0298:phyxminiproblems-problemphyxmini0298-figurepoint}
  \lean{PhyXMiniProblems.ProblemPhyXMini0298.FigurePoint}
  The two labelled endpoints of the horizontal vibrating segment
\end{definition}
\begin{proof}
  This is the declared model definition of Figure Point.
\end{proof}
\begin{definition}[Figure Component]
  \label{def:physics:phyx-mini-0298:phyxminiproblems-problemphyxmini0298-figurecomponent}
  \lean{PhyXMiniProblems.ProblemPhyXMini0298.FigureComponent}
  Individually identifiable objects and portions of string in the figure
\end{definition}
\begin{proof}
  This is the declared model definition of Figure Component.
\end{proof}
\begin{definition}[String Segment]
  \label{def:physics:phyx-mini-0298:phyxminiproblems-problemphyxmini0298-stringsegment}
  \lean{PhyXMiniProblems.ProblemPhyXMini0298.StringSegment}
  The two string portions separated by the support at Q
\end{definition}
\begin{proof}
  This is the declared model definition of String Segment.
\end{proof}
\begin{definition}[String Route]
  \label{def:physics:phyx-mini-0298:phyxminiproblems-problemphyxmini0298-stringroute}
  \lean{PhyXMiniProblems.ProblemPhyXMini0298.StringRoute}
  The complete route of the single string drawn in the primary figure
\end{definition}
\begin{proof}
  This is the declared model definition of String Route.
\end{proof}
\begin{definition}[Oscillator String Setup]
  \label{def:physics:phyx-mini-0298:phyxminiproblems-problemphyxmini0298-oscillatorstringsetup}
  \lean{PhyXMiniProblems.ProblemPhyXMini0298.OscillatorStringSetup}
  \uses{def:physics:phyx-mini-0298:phyxminiproblems-problemphyxmini0298-massquantity, def:physics:phyx-mini-0298:phyxminiproblems-problemphyxmini0298-lengthquantity, def:physics:phyx-mini-0298:phyxminiproblems-problemphyxmini0298-frequencyquantity, def:physics:phyx-mini-0298:phyxminiproblems-problemphyxmini0298-accelerationquantity, def:physics:phyx-mini-0298:phyxminiproblems-problemphyxmini0298-tensionquantity, def:physics:phyx-mini-0298:phyxminiproblems-problemphyxmini0298-linearmassdensityquantity, def:physics:phyx-mini-0298:phyxminiproblems-problemphyxmini0298-speedquantity, def:physics:phyx-mini-0298:phyxminiproblems-problemphyxmini0298-figurepoint, def:physics:phyx-mini-0298:phyxminiproblems-problemphyxmini0298-figurecomponent, def:physics:phyx-mini-0298:phyxminiproblems-problemphyxmini0298-stringsegment, def:physics:phyx-mini-0298:phyxminiproblems-problemphyxmini0298-stringroute}
  Independent physical quantities and qualitative data for the pictured apparatus. In particular, hangingMass is not assigned the requested numerical value here
\end{definition}
\begin{proof}
  This is the declared model definition of Oscillator String Setup.
\end{proof}
\begin{definition}[Matches Problem And Primary Figure]
  \label{def:physics:phyx-mini-0298:phyxminiproblems-problemphyxmini0298-matchesproblemandprimaryfigure}
  \lean{PhyXMiniProblems.ProblemPhyXMini0298.MatchesProblemAndPrimaryFigure}
  \uses{def:physics:phyx-mini-0298:phyxminiproblems-problemphyxmini0298-lengthinmeters, def:physics:phyx-mini-0298:phyxminiproblems-problemphyxmini0298-frequencyinhertz, def:physics:phyx-mini-0298:phyxminiproblems-problemphyxmini0298-linearmassdensityinkilogramspermeter, def:physics:phyx-mini-0298:phyxminiproblems-problemphyxmini0298-oscillatorstringsetup}
  The numerical data and categorical facts stated in the problem or visible in the primary figure. The density readout 1 / 625 kg/m is the SI conversion of 1.6 g/m. The requested mass and all four answer values are absent
\end{definition}
\begin{proof}
  This is the declared model definition of Matches Problem And Primary Figure.
\end{proof}
\begin{definition}[Uses Standard Earth Gravity]
  \label{def:physics:phyx-mini-0298:phyxminiproblems-problemphyxmini0298-usesstandardearthgravity}
  \lean{PhyXMiniProblems.ProblemPhyXMini0298.UsesStandardEarthGravity}
  \uses{def:physics:phyx-mini-0298:phyxminiproblems-problemphyxmini0298-accelerationinmeterspersecondsquared, def:physics:phyx-mini-0298:phyxminiproblems-problemphyxmini0298-oscillatorstringsetup}
  The conventional near-Earth gravitational acceleration used by the choices
\end{definition}
\begin{proof}
  This is the declared model definition of Uses Standard Earth Gravity.
\end{proof}
\begin{definition}[Has Positive Physical Parameters]
  \label{def:physics:phyx-mini-0298:phyxminiproblems-problemphyxmini0298-haspositivephysicalparameters}
  \lean{PhyXMiniProblems.ProblemPhyXMini0298.HasPositivePhysicalParameters}
  \uses{def:physics:phyx-mini-0298:phyxminiproblems-problemphyxmini0298-massinkilograms, def:physics:phyx-mini-0298:phyxminiproblems-problemphyxmini0298-lengthinmeters, def:physics:phyx-mini-0298:phyxminiproblems-problemphyxmini0298-frequencyinhertz, def:physics:phyx-mini-0298:phyxminiproblems-problemphyxmini0298-accelerationinmeterspersecondsquared, def:physics:phyx-mini-0298:phyxminiproblems-problemphyxmini0298-tensioninnewtons, def:physics:phyx-mini-0298:phyxminiproblems-problemphyxmini0298-linearmassdensityinkilogramspermeter, def:physics:phyx-mini-0298:phyxminiproblems-problemphyxmini0298-speedinmeterspersecond, def:physics:phyx-mini-0298:phyxminiproblems-problemphyxmini0298-oscillatorstringsetup}
  Positivity and nondegeneracy conditions for the taut-string model
\end{definition}
\begin{proof}
  This is the declared model definition of Has Positive Physical Parameters.
\end{proof}
\begin{definition}[Satisfies Ideal Support String And Harmonic Laws]
  \label{def:physics:phyx-mini-0298:phyxminiproblems-problemphyxmini0298-satisfiesidealsupportstringandharmoniclaws}
  \lean{PhyXMiniProblems.ProblemPhyXMini0298.SatisfiesIdealSupportStringAndHarmonicLaws}
  \uses{def:physics:phyx-mini-0298:phyxminiproblems-problemphyxmini0298-massinkilograms, def:physics:phyx-mini-0298:phyxminiproblems-problemphyxmini0298-lengthinmeters, def:physics:phyx-mini-0298:phyxminiproblems-problemphyxmini0298-frequencyinhertz, def:physics:phyx-mini-0298:phyxminiproblems-problemphyxmini0298-accelerationinmeterspersecondsquared, def:physics:phyx-mini-0298:phyxminiproblems-problemphyxmini0298-tensioninnewtons, def:physics:phyx-mini-0298:phyxminiproblems-problemphyxmini0298-linearmassdensityinkilogramspermeter, def:physics:phyx-mini-0298:phyxminiproblems-problemphyxmini0298-speedinmeterspersecond, def:physics:phyx-mini-0298:phyxminiproblems-problemphyxmini0298-oscillatorstringsetup}
  The standard textbook idealizations used in the calculation: a massless string over a frictionless support transmits the same tension; static balance makes the vertical tension equal to the hanging weight; a taut string obeys v = sqrt (T / μ); a length-L segment with nodes at both ends has fₙ = n v / (2 L) in harmonic n. These laws are stated generically and contain neither the requested mass nor the recorded answer 0.846 kg
\end{definition}
\begin{proof}
  This is the declared model definition of Satisfies Ideal Support String And Harmonic Laws.
\end{proof}
\begin{lemma}[fourth Harmonic Wave Speed exact]
  \label{lem:physics:phyx-mini-0298:phyxminiproblems-problemphyxmini0298-fourthharmonicwavespeed-exact}
  \lean{PhyXMiniProblems.ProblemPhyXMini0298.fourthHarmonicWaveSpeed_exact}
  \uses{def:physics:phyx-mini-0298:phyxminiproblems-problemphyxmini0298-speedinmeterspersecond, def:physics:phyx-mini-0298:phyxminiproblems-problemphyxmini0298-oscillatorstringsetup, def:physics:phyx-mini-0298:phyxminiproblems-problemphyxmini0298-matchesproblemandprimaryfigure, def:physics:phyx-mini-0298:phyxminiproblems-problemphyxmini0298-haspositivephysicalparameters, def:physics:phyx-mini-0298:phyxminiproblems-problemphyxmini0298-satisfiesidealsupportstringandharmoniclaws}
  The fourth-harmonic relation with L = 1.20 m and f = 120 Hz gives the unrounded transverse wave speed v = 72 m/s
\end{lemma}
\begin{proof}
  The conclusion follows from the governing relations and figure readouts named in the statement of fourth Harmonic Wave Speed exact.
\end{proof}
\begin{lemma}[required Hanging Mass exact]
  \label{lem:physics:phyx-mini-0298:phyxminiproblems-problemphyxmini0298-requiredhangingmass-exact}
  \lean{PhyXMiniProblems.ProblemPhyXMini0298.requiredHangingMass_exact}
  \uses{def:physics:phyx-mini-0298:phyxminiproblems-problemphyxmini0298-massinkilograms, def:physics:phyx-mini-0298:phyxminiproblems-problemphyxmini0298-oscillatorstringsetup, def:physics:phyx-mini-0298:phyxminiproblems-problemphyxmini0298-matchesproblemandprimaryfigure, def:physics:phyx-mini-0298:phyxminiproblems-problemphyxmini0298-usesstandardearthgravity, def:physics:phyx-mini-0298:phyxminiproblems-problemphyxmini0298-haspositivephysicalparameters, def:physics:phyx-mini-0298:phyxminiproblems-problemphyxmini0298-satisfiesidealsupportstringandharmoniclaws}
  Using μ = 0.0016 kg/m, v = 72 m/s, and g = 9.8 m/s², the required unrounded hanging mass is 5184 / 6125 kg
\end{lemma}
\begin{proof}
  The conclusion follows from the governing relations and figure readouts named in the statement of required Hanging Mass exact.
\end{proof}
\begin{definition}[Answer Choice]
  \label{def:physics:phyx-mini-0298:phyxminiproblems-problemphyxmini0298-answerchoice}
  \lean{PhyXMiniProblems.ProblemPhyXMini0298.AnswerChoice}
  Labels printed beside the four multiple-choice answers
\end{definition}
\begin{proof}
  This is the declared model definition of Answer Choice.
\end{proof}
\begin{definition}[mass In Kilograms]
  \label{def:physics:phyx-mini-0298:phyxminiproblems-problemphyxmini0298-answerchoice-massinkilograms}
  \lean{PhyXMiniProblems.ProblemPhyXMini0298.AnswerChoice.massInKilograms}
  \uses{def:physics:phyx-mini-0298:phyxminiproblems-problemphyxmini0298-massinkilograms, def:physics:phyx-mini-0298:phyxminiproblems-problemphyxmini0298-answerchoice}
  Displayed mass beside each answer label, in kilograms
\end{definition}
\begin{proof}
  This is the declared model definition of mass In Kilograms.
\end{proof}
\begin{definition}[recorded Answer Choice]
  \label{def:physics:phyx-mini-0298:phyxminiproblems-problemphyxmini0298-recordedanswerchoice}
  \lean{PhyXMiniProblems.ProblemPhyXMini0298.recordedAnswerChoice}
  \uses{def:physics:phyx-mini-0298:phyxminiproblems-problemphyxmini0298-answerchoice}
  The answer label recorded in the supplied dataset metadata
\end{definition}
\begin{proof}
  This is the declared model definition of recorded Answer Choice.
\end{proof}
\begin{definition}[Matches Answer Choice]
  \label{def:physics:phyx-mini-0298:phyxminiproblems-problemphyxmini0298-matchesanswerchoice}
  \lean{PhyXMiniProblems.ProblemPhyXMini0298.MatchesAnswerChoice}
  \uses{def:physics:phyx-mini-0298:phyxminiproblems-problemphyxmini0298-massinkilograms, def:physics:phyx-mini-0298:phyxminiproblems-problemphyxmini0298-oscillatorstringsetup, def:physics:phyx-mini-0298:phyxminiproblems-problemphyxmini0298-answerchoice}
  Agreement with a mass displayed to the nearest 0.001 kg. The tolerance is 0.0005 kg, half a unit in the last displayed digit
\end{definition}
\begin{proof}
  This is the declared model definition of Matches Answer Choice.
\end{proof}
% --- Archon declaration topology end ---
% --- Archon physics formalization source end ---
